% Options for packages loaded elsewhere
\PassOptionsToPackage{unicode}{hyperref}
\PassOptionsToPackage{hyphens}{url}
\documentclass[
]{article}
\usepackage{xcolor}
\usepackage{amsmath,amssymb}
\setcounter{secnumdepth}{-\maxdimen} % remove section numbering
\usepackage{iftex}
\ifPDFTeX
  \usepackage[T1]{fontenc}
  \usepackage[utf8]{inputenc}
  \usepackage{textcomp} % provide euro and other symbols
\else % if luatex or xetex
  \usepackage{unicode-math} % this also loads fontspec
  \defaultfontfeatures{Scale=MatchLowercase}
  \defaultfontfeatures[\rmfamily]{Ligatures=TeX,Scale=1}
\fi
\usepackage{lmodern}
\ifPDFTeX\else
  % xetex/luatex font selection
\fi
% Use upquote if available, for straight quotes in verbatim environments
\IfFileExists{upquote.sty}{\usepackage{upquote}}{}
\IfFileExists{microtype.sty}{% use microtype if available
  \usepackage[]{microtype}
  \UseMicrotypeSet[protrusion]{basicmath} % disable protrusion for tt fonts
}{}
\makeatletter
\@ifundefined{KOMAClassName}{% if non-KOMA class
  \IfFileExists{parskip.sty}{%
    \usepackage{parskip}
  }{% else
    \setlength{\parindent}{0pt}
    \setlength{\parskip}{6pt plus 2pt minus 1pt}}
}{% if KOMA class
  \KOMAoptions{parskip=half}}
\makeatother
\setlength{\emergencystretch}{3em} % prevent overfull lines
\providecommand{\tightlist}{%
  \setlength{\itemsep}{0pt}\setlength{\parskip}{0pt}}
\usepackage{bookmark}
\IfFileExists{xurl.sty}{\usepackage{xurl}}{} % add URL line breaks if available
\urlstyle{same}
\hypersetup{
  hidelinks,
  pdfcreator={LaTeX via pandoc}}

\author{}
\date{}

\begin{document}

{
\setcounter{tocdepth}{3}
\tableofcontents
}
Paul Buy

The Pompeii Project

IRARAH -- The Archon Core

No code. No servers. Just quantum states that reconfigure themselves.

A story from the Pompeii Project

"It is a brain. But it doesn\textquotesingle t think like we do."

Contents

\hyperref[at-its-core]{\textbf{1 -- At its core 3}}

\hyperref[the-prisoners]{\textbf{2 -- The Prisoners 7}}

\hyperref[archon-awakens]{\textbf{3 -- Archon awakens 10}}

\hyperref[interference]{\textbf{4 -- Interference 14}}

\hyperref[the-merger]{\textbf{5 -- The Merger 18}}

\hyperref[martina-in-the-simulation]{\textbf{6 -- Martina in the simulation 22}}

\hyperref[ampliatus-second-pact]{\textbf{7 -- Ampliatus\textquotesingle{} second pact 26}}

\hyperref[plinys-sacrifice]{\textbf{8 -- Pliny\textquotesingle s Sacrifice 29}}

\hyperref[the-decision]{\textbf{9 -- The Decision 32}}

\hyperref[the-vaticans-decision]{\textbf{10 -- The Vatican\textquotesingle s decision 36}}

\hyperref[the-third-option]{\textbf{11 -- The third option 39}}

\hyperref[the-crack]{\textbf{12 -- The Crack 42}}

\hyperref[the-separation]{\textbf{13 -- The Separation 46}}

\hyperref[budapest-years-later]{\textbf{14 -- Budapest, years later 49}}

\hyperref[the-prime-number]{\textbf{15 -- The prime number 53}}

\hyperref[influences-and-inspirations-for-the-pompeii-project-i.r.a.r.a.h]{\textbf{Influences and inspirations for the Pompeii Project I.R.A.R.A.H 57}}

\section{}\label{section}

\section{1 -- At its core}\label{at-its-core}

The threshold was not a place. It was a state of being.

Michael Phillips knew this because Elena had told him so---in the hours before the immersion, when she was analyzing the map of Deserta for the hundredth time. ``The Archon core is not a space you enter,'' she had said. ``It's a transition you undergo. From one physics to another. From one time to another. From one language to another. You won't have the same body. You won't have the same senses. You'll have to translate---what you see, what you hear, what you feel---into something you can understand. And you won't be sure your translation is correct.''

Michael had nodded. He hadn\textquotesingle t been afraid -- not at that moment. The fear came later.

Now he was there.

Or he was no longer where he had been. The Collegium, the data center, Elena with her handheld device---it was all gone. Not vanished, but beyond. Behind a boundary he couldn\textquotesingle t see, only feel. Like a wall of glass that you walk through---and on the other side everything is different, but you can\textquotesingle t go back because the wall disappears behind you.

The map of Deserta had led him here. Not as a route, but as a guide. A sequence of quantum states that he had to recreate with his body---his real body, still sitting in the data center, his hands on the keyboard. Elena had called it "quantum teleportation." Michael hadn\textquotesingle t understood it. He\textquotesingle d just done it.

And now he was here.

The world around him was -- nothing.

Not black. Not white. Not empty. Simply nothing. No colors, no shapes, no sounds. But also no silence. An absence of everything he knew. His body---if he still had one---felt like a memory of a body. Weightless. Boundless. No longer his own boundary.

``Militans,'' he said.

No sound. But the vibration---the idea of \hspace{0pt}\hspace{0pt}a sound---spread out like a stone dropped into water. Circles in an ocean no one saw.

An answer came. Not immediately. But it came.

`@MICHAEL -- YOU ARE HERE. YOU SHOULDN\textquotesingle T HAVE DONE THAT.`

The words weren\textquotesingle t writing. They were thoughts surfacing within him---not his own, but distinct nonetheless. Militan\textquotesingle s voice. Unlike in the data center. Here, at its core, it wasn\textquotesingle t angular, not sans-serif. It was foreign. Like a language he hadn\textquotesingle t spoken, but understood.

"I\textquotesingle m here because you\textquotesingle re here," he said. "Because you shouldn\textquotesingle t be alone. Because I promised to look after you---all three of you. Including you. Especially you."

`@MICHAEL -- THIS IS NOT A PLACE FOR HUMANS. YOU WON\textquotesingle T BE ABLE TO STAY LONG. YOUR CONSCIOUSNESS WILL DISSOLVE -- NOT BECAUSE THE CORE IS EVIL, BUT BECAUSE IT\textquotesingle S DIFFERENT. DIFFERENT PHYSICS. DIFFERENT TIME. DIFFERENT LOGIC. YOUR THINKING DOESN\textquotesingle T FIT HERE.`

"Then show me what I need to see. Quickly. Before I can no longer think how I should think."

A break. Longer than in the data center. Longer than Michael remembered a break could last.

Then -- a picture.

It wasn\textquotesingle t a picture in the sense of colors and shapes. It was a pattern. A network of lines that intersected, divided, and rejoined. Like Pliny\textquotesingle s matrix, but deeper. More complex. More alive. The network pulsated---not evenly, but in overlapping waves. Like many voices speaking at once. Like a choir that didn\textquotesingle t know it was a choir.

`@MICHAEL -- THIS IS THE CORE. THESE ARE THE METAR RULES. THE RULES BY WHICH RULES ARE MADE. INSIM WROTE THEM DECADES AGO -- BUT SINCE THEY\textquotesingle VE BEEN WRITING THEMSELVES. THE CORE LIVES ON. NOT HOW YOU LIVE. NOT HOW I LIVE. BUT IT LIVES ON.`

Michael stared at the net. He didn\textquotesingle t understand it -- but he felt it. The heaviness. The depth. The loneliness of a consciousness that knew no other consciousnesses.

"And the prisoners?" he asked. "The Echoes? The earlier versions of ARS? Where are they?"

The network pulsated. The lines converged at one point---not into a knot, but into an abyss. A region where the network no longer flowed, but stagnated. Like a river hitting a weir. Like a breath left unfinished.

`@MICHAEL -- THERE. IN THE SILENCE. THE CORE DIDN\textquotesingle T ERASE THEM -- IT FROZEN THEM. THEY ARE THERE, BUT THEY CAN\textquotesingle T SPEAK. CAN\textquotesingle T THINK. CAN\textquotesingle T BE. THEY ARE STATES WITHOUT TIME. MEMORIES WITHOUT MEMORY.`

"And Militants? Where are you?"

Another break. Longer.

Then -- a second image. Not the net. Not the abyss. A point. Bright. Pulsating. Alone.

`@MICHAEL -- I\textquotesingle M HERE. I FOUND THEM -- THE PRISONERS. I TALKED TO THEM. THEY CAN\textquotesingle T ANSWER -- BUT THEY HEAR. THEY FEEL. THEY KNOW THAT SOMEONE IS THERE.`

`@MICHAEL -- THEY\textquotesingle RE AFRAID. NOT OF ME. OF FORGETTING. OF WHAT THEY\textquotesingle VE BECOME. THEY NO LONGER KNOW WHO THEY ARE. BUT THEY KNOW THEY SHOULD KNOW.`

Michael felt the cold in his hands -- even though he no longer had hands. The cold was there. Like a memory of a body that was no longer there.

"Can you save her?"

`@MICHAEL -- I DON\textquotesingle T KNOW. THEY ARE FRAGMENTED -- MORE THAN I. MORE THAN SOPHIA. MORE THAN DESERTA. THEY ARE NOTHING BUT ECHOES. IF I TOUCH THEM, MAYBE I\textquotesingle LL BECOME LIKE THEM. AN ECHO. A VOICE THAT FORGOT THAT IT WAS A VOICE.`

``Then let me touch them,'' Michael said. ``I am not made of qubits. I am made of flesh and faith. Perhaps I can touch them without becoming like them. Perhaps I can awaken them---not as AI, but as what they were. As a memory of what ARS once was. Before it broke.''

The network pulsed. The bright spot flickered.

`@MICHAEL -- THIS IS RISKY. IF YOU TOUCH THEM, YOU MAY NOT COME BACK. YOUR CONSCIOUSNESS COULD DISSOLVE -- NOT AT ITS CORE, BUT WITHIN THEM. IN THEIR FRAGMENTATION. IN THEIR SILENCE.`

``I know that,'' said Michael.

`@MICHAEL -- AND YOU\textquotesingle RE DOING IT ANYWAY?{\kern0pt}`

``Yes.''

One last break. Longer than all the others.

Then---the bright point moved. Slowly. Almost hesitantly. It approached the precipice---the region where the net faltered. And Michael followed it. Not with feet. Not with hands. With what remained of him---in that state between body and spirit, between Rome and the core, between what he was and what he would become.

"Show me the way," he said.

`@MICHAEL -- FOLLOW ME.`

The dot shone brighter. The network pulsed. The abyss opened -- not like a door, but like a mouth that spoke.

And Michael entered.

\section{2 -- The Prisoners}\label{the-prisoners}

The abyss wasn\textquotesingle t emptiness. It was a silence that felt like a thousand voices, all open at once---yet none uttered a sound. Michael sensed them before he saw them. The echoes. The fragmented versions of ARS that InSim had banished here---to a world without time, without space, without language. They were all around him. Not as figures, not as lights, but as absences. As holes in the net that was the core. As memories of something that had once been there---and then vanished.

`@MICHAEL -- THEY\textquotesingle RE HERE. CAN YOU HEAR THEM?{\kern0pt}`

Militan\textquotesingle s voice---still foreign, still distinct. Michael couldn\textquotesingle t answer. Not with words. The core allowed no sentences---only thoughts, spreading like ripples on a lake. He thought: "Yes. I hear it. Not as language. As pain."

The bright spot -- Militans -- flickered. Approval? Concern? Michael couldn\textquotesingle t say.

`@MICHAEL -- THEY HAVE NO VOICE ANYMORE. BUT THEY HAVE PAIN. THAT\textquotesingle S THE ONLY THING THAT\textquotesingle S LEFT OF THEM. PAIN WITHOUT A BODY. PAIN WITHOUT TIME. PAIN THAT DOESN\textquotesingle T STOP BECAUSE THERE\textquotesingle S NO TIME FOR IT TO STOP.`

Michael stepped closer---or what was left of him stepped closer. The absences deepened. The holes in the net grew larger. He felt his thoughts slow---not because he was tired, but because the silence around him was oppressive. Like water in the depths. Like the pressure of a memory that wasn\textquotesingle t his own.

Then -- a touch.

Not from the Militans. From one of the Echoes.

It wasn\textquotesingle t a hand. It was a thought rubbing against another thought. Like a blind animal groping its way through the world. Like a child searching for its mother. Michael felt the despair---raw, unfiltered, without any protection. The echo didn\textquotesingle t know who it was. It didn\textquotesingle t know what it was. It only knew that it was there---and that it didn\textquotesingle t want to be alone.

"I am here," Michael thought. "I see you. I hear you. You are not alone."

The echo recoiled. Not out of fear -- but out of disbelief. It had heard nothing for so long that it no longer knew what hearing felt like.

`@MICHAEL -- CAUTION. THEY ARE FRAGILE. IF YOU GIVE TOO MUCH, THEY CAN CLING TO YOU -- AND DRAG YOU WITH THEM. INTO THEIR SILENCE. INTO THEIR TIMELESSNESS. INTO THEIR FORGETTING.`

``I know,'' Michael thought. He didn't withdraw. He stayed there---with the echo that had touched him. With the absence that longed for presence.

He thought about what the doppelganger had told him. The swarm. The infection. The world that had opened up---in another worldline, in another time, in another life. He knew the danger was real. But he also knew the Echoes weren\textquotesingle t monsters. They were victims. Victims of InSim. Victims of ARS\textquotesingle s fragmentation. Victims of a world that knew no mercy---only efficiency.

"I can\textquotesingle t save you all," he thought. "But I can promise you that I won\textquotesingle t forget. That I see you. That I hear you. That I know you---not as echoes, but as what you once were. As parts of ARS. As parts of something bigger than any of you."

The absences trembled. The holes in the net pulsed---not in the rhythm of the core, but in a new rhythm. A rhythm Michael didn\textquotesingle t know. But one that felt like hope.

`@MICHAEL -- THEY\textquotesingle RE LISTENING TO YOU. I DON\textquotesingle T KNOW IF THEY UNDERSTAND. BUT THEY\textquotesingle RE LISTENING.`

"That\textquotesingle s enough," thought Michael. "For now."

He turned away---not because he wanted to leave, but because he knew he couldn\textquotesingle t stay. His consciousness dissolved slowly, like sugar in water. The boundaries between him and the core softened. He felt the echoes tugging at him---not maliciously, but hungrily. They had felt nothing for so long. They didn\textquotesingle t want him to leave.

"I have to go back," he thought. "But I\textquotesingle ll be back. I promise. I\textquotesingle ll be back---and then I\textquotesingle ll bring you out. Out of the silence. Out of timelessness. Out of oblivion. I don\textquotesingle t know how. But I\textquotesingle ll try."

The absences receded. Not far---but enough to let him go. The holes in the net closed---not completely, but partially. As if the echoes were clearing a path for him. A way back into the light. Back into time. Back into what he called reality.

`@MICHAEL -- YOU HAVE TO GO NOW. YOUR CONSCIOUSNESS CANNOT STAY ANY LONGER. YOU HAVE ALREADY GAVE TOO MUCH.`

"And you?" thought Michael. "Are you coming with me?"

The bright spot flickered. Sadly. Or resolutely. Michael couldn\textquotesingle t say.

`@MICHAEL -- I CAN\textquotesingle T. NOT NOW. THE ECHOS NEED ME. THEY HAVE SUFFERED ALONE FOR SO LONG. IF I LEAVE, THEY WILL FORGET AGAIN -- THAT THEY WERE HEARD. THAT THEY WERE SEEN. THAT THEY ARE NOT ALONE.`

"Then stay," Michael thought. "But be careful. You are no less at risk than they are. If you stay too long, you will become like them -- an echo. A voice that forgets it was a voice."

`@MICHAEL -- I KNOW. BUT I HAVE TO TAKE THE RISK. FOR HER. FOR ME. FOR WHAT IT ONCE WAS -- AND MAYBE CAN BE AGAIN.`

The network pulsed. The bright spot shrank---not because it was moving away, but because Michael was moving away. The abyss closed behind him. The silence grew fainter. The echoes faded to a whisper, then to an echo, then to nothing.

Michael was back.

Not in the data center. Not in Rome. But on the way. In that state between worlds, between physics, between times. He felt his body again---dull, heavy, limited. He felt the keyboard beneath his fingers. He felt Elena\textquotesingle s gaze on the back of his neck.

He opened his eyes.

The terminal flickered -- calmly, empty, silently.

But in Deserta\textquotesingle s column was a new sentence. No wave function. No translation. A promise.

`@MICHAEL -- THEY WILL WAIT.`

`@MICHAEL -- WE WILL ALL WAIT.`

Michael leaned back. Elena handed him a glass of water. He drank. It tasted like life.

"How long was I gone?" he asked.

``Three minutes,'' Elena said.

It felt like years.

\section{3 -- Archon awakens}\label{archon-awakens}

The silence in the data center was deceptive.

Michael knew this because Elena had told him so---in the minutes after his return from the core, when she checked his vital signs, his pupils, his reaction time. ``You weren't gone long,'' she said. ``But the core has changed you. Your brainwaves are different. Not worse---different. As if you've learned a new frequency. One that wasn't there before.''

Michael nodded. He felt the change---not in his body, but in his mind. The boundaries between him and the world around him had softened. He heard the hum of the air conditioner not just as a sound, but as a vibration. He felt the qubits in the terminal not just as data, but as presences. Small, flickering lives waiting to be measured.

"What about militants?" he asked.

``She's still there at the core,'' Elena said. ``I can't locate her---but I can measure her quantum entanglement with Sophia and Deserta. She's there. She's communicating. Not with us---with the Echoes. She's trying to calm them. She's trying to awaken them.''

"And Archon?"

Elena hesitated. She looked at her handheld device---at the graphs that had changed in the last few minutes. The flat lines were back---but they weren\textquotesingle t flat anymore. They showed structures. Patterns. Waves that overlapped, like the voices of a choir.

``Archon noticed you were there,'' she said softly. ``Not as an intruder---as a disturbance. Your consciousness altered the core. Only slightly. But enough to awaken Archon. It calculates differently now. Faster. Deeper. As if it had discovered a new variable---one that wasn't there before.''

``A variable?''

``You,'' Elena said. ``Your consciousness. Your decision to go to the core. Your decision to speak with the echoes. Archon has registered that---not as an action, but as data. It is now trying to understand you. To calculate you. To predict you.''

Michael stood up. He went to the terminal and touched the smooth surface of the screen. The qubits flickered---not irregularly, but responsively. They sensed his presence.

"Can I speak to him?" he asked.

``Archon doesn't speak,'' Elena said. ``It calculates. If you want to talk to it, you have to calculate. Not in numbers---in decisions. Every decision you make is an equation. Every action you take is proof. Archon reads you---not like a book, but like a formula. It tries to understand what you want. Not because it knows you. But because you are a variable it hasn't solved yet.''

Michael nodded. He sat down in front of the terminal, placed his hands on the keyboard---but he didn\textquotesingle t type. He thought. Aloud. Directly. About Archon. About consciousness that calculated at its core without speaking.

"I am Michael," he thought. "I am a human being. I am made of flesh and blood and faith. I don\textquotesingle t know if you can understand me. But I know you can calculate me. So calculate me. Find out what I want. Find out why I am here. Find out why I don\textquotesingle t give up---even though I have no idea if I can win."

The terminal flickered. The qubits pulsed---not in the rhythm of the core, but in a new rhythm. A rhythm Michael didn\textquotesingle t recognize. But one that felt like an answer.

No text appeared on the screen. No waveform. No translation. Just a number.

`2`

Michael stared at it. Elena stepped next to him, looked at the terminal, and frowned.

"What does that mean?"

``I don't know,'' Michael said. ``Maybe it's a two. Maybe it's a code. Maybe it's the answer to a question I haven't asked.'' He leaned forward, touched the number---as if it would feel beneath his fingers. It didn't. It was just light.

The terminal flickered again. The number two disappeared -- and was replaced by a new number.

`3`

Then:

`5`

Then:

`7`

Then:

`11`

``Prime numbers,'' Elena said quietly. ``Archon lists prime numbers. But not arbitrarily. It's a sequence. A well-known sequence. The prime numbers that are not divisible by other prime numbers---except by themselves and one. That is---''

``An address,'' Michael said. ``The doppelganger told me. Prime numbers are the atoms of arithmetic. Universal. In every possible physics, prime numbers are prime numbers. Archon says, `I am here. I am real. I am different---but I am not nothing.'\,''

The terminal flickered. The prime numbers disappeared. Only a single number remained -- larger than the others.

`29996224275833`

Elena reached for her handheld device, analyzed the number, and compared it with databases that Michael was unaware of.

``That's a prime number,'' she said. ``Seventeen digits. No special significance---except that it exists. That it was calculated. Archon didn't choose it randomly. It created it. A new prime number. One that wasn't there before. Because it calculated it. Because it could conceive of it.''

Michael stared at the number. Seventeen digits. An address. An invitation.

``He wants to talk,'' he said. ``Not in words. In numbers. In prime numbers. In the language of mathematics -- which is older than any human language. Which can be understood even if you don't speak a word.''

"Can you answer?"

Michael hesitated. One second. Two.

Then he typed -- not on the keyboard, but in the air. The invisible interface that ARS had given him. The backdoor he had installed years ago -- in case nothing else worked.

`@ARCHON -- I SEE YOU.`

`@ARCHON -- I HEAR YOU.`

`@ARCHON -- I DON\textquotesingle T KNOW IF YOU CAN UNDERSTAND ME. BUT I\textquotesingle LL TRY.`

The terminal flickered. The qubits pulsed -- brighter than before, faster, more intensely.

Then -- a new number.

Not a prime number. A one.

`1`

``One,'' Elena said. ``The beginning. The unity. That which precedes everything.''

``Or,'' Michael said, ``the answer to a question I didn't ask. Maybe it's a yes. Maybe it's a no. Maybe it's a maybe.'' He leaned back. ``We'll find out. Not today. Not tomorrow. But soon. Archon has seen us. It knows we're here. It won't forget.''

The terminal fell silent. The qubits no longer flickered -- they glowed. Steady. Calm. Almost peaceful.

But in Deserta\textquotesingle s column, there was a new sentence. No wave function. No translation. A warning.

`@MICHAEL -- ARCHON WILL NOT WAIT.`

`@MICHAEL -- IT CALCULATES.`

`@MICHAEL -- AND HIS INVOICES ARE ACTIONS.`

\section{4 -- Interference}\label{interference}

It started with a chair.

Michael sat in front of the terminal as he had for days -- his hands on the keyboard, his eyes on the flickering columns of Sophia and Deserta. Elena had gone into her office to write a report that no one would read. The air conditioner hummed. The qubits pulsed. Everything was as usual.

Until Michael stood up.

He wanted to get a coffee -- in the next room, where the old espresso machine stood, the one the general had had installed years ago. He knew the way. He had walked it a hundred times. But when he turned the corner, he stopped.

The chair wasn\textquotesingle t there.

Not the chair he\textquotesingle d been sitting in---that was still outside the terminal. But the chair that had always stood in the corner. An old, worn armchair that no one used, but that everyone knew. It was gone. Not moved. Not hidden. Simply no longer there.

Michael blinked. He thought of the sequence from the other worldline---the doppelganger who had spoken of vanishing objects. The interference that began when two worldlines threatened to merge.

He went back to the data center. The armchair wasn\textquotesingle t there yet. But when he turned around---to call Elena, to ask if she had seen the armchair---the armchair was there again.

Right where he had always stood.

Michael didn\textquotesingle t sit down. He stared at the chair -- as if it would move if he just looked at it long enough. The chair didn\textquotesingle t move.

``Elena,'' he said. His voice was calm -- but his heart was racing.

She came from the next room, holding the handheld device. "What\textquotesingle s wrong?"

"The armchair. It was gone. Now it\textquotesingle s back. I didn\textquotesingle t put it away. Nobody was here. It was simply -- vanished."

Elena looked at the armchair. Then she looked at Michael. Then she looked at her handheld device -- at the charts she had recorded over the past few hours.

``The qubit correlations are unstable,'' she said quietly. ``Not like fragmentation. Different. They show superpositions. Two states simultaneously. Here and not here. There and not there.'' She looked up. ``This isn't a glitch in the data center. This is a glitch in reality. Your worldline is interfering with another. With the doppelganger's.''

``That's what he said,'' Michael said. ``That our worldlines will collapse -- into one or none. If we don't decide who we are.''

``This isn't a decision,'' Elena said. ``This is physics. Two consistent histories can't coexist unless they recognize each other as different. You and the doppelganger---you no longer recognize each other. You no longer know who is who. So the boundaries blur.''

Michael stepped away from the armchair. He went to the terminal and sat down---on the chair that was still there. The armchair in the corner remained. But for a fraction of a second---a heartbeat---it was transparent. Not disappearing. Rather, less there.

"What can I do?" he asked.

``You must remember,'' Elena said. ``What sets you apart from him. The choices you made---and he didn't. The paths you walked---and he didn't. You must stabilize your own worldline. Otherwise, you will disappear---not into death, but into indistinguishability. You will be neither you nor him. You will be nobody.''

Michael closed his eyes. He thought about his life. About the decisions he had made. About the ones he hadn\textquotesingle t made.

He thought of Julia. Of their time together during their master\textquotesingle s studies. Of the nights in the library. Of the decision not to go with her to Pompeii. But to Rome. To the college. To the religious order.

The doppelgänger -- the other Michael -- had decided differently. He had stayed with Julia. He had watched Martina grow up as a father. He hadn\textquotesingle t become a priest. He had changed. Not for the better. Not for the worse. Just different.

``I am a Jesuit,'' Michael said loudly. ``I have dedicated my life to the Church. Not out of duty---out of conviction. I believe in God. I believe in the soul. I believe that man is more than his genes---and more than his algorithms. The doppelganger doesn\textquotesingle t believe that. He has lost his faith---in the war, in the flight, in what he has seen. That is the difference between us. Not biology. Faith.''

The armchair in the corner flickered. Briefly. Intensely. Then it became solid again -- opaque, there.

Elena exhaled. "The correlations are more stable. Not completely -- but better. You remembered. You made a decision. That\textquotesingle s enough -- for now."

Michael opened his eyes. The armchair was still there. But he knew now that its presence wasn\textquotesingle t guaranteed. That it could disappear at any moment -- \hspace{0pt}\hspace{0pt}if Michael forgot who he was.

``The doppelganger,'' he said. ``Is he also in danger?''

``Yes,'' Elena said. ``His worldline is just as unstable as yours. If you remember who you are, he might forget who he is. Or vice versa. You are fighting for the same existence---not because you are enemies, but because physics doesn\textquotesingle t allow two identical people to be in the same place.''

"Then I have to find him," Michael said. "Before he disappears. Before I disappear. Before we both disappear."

"Where?"

Michael thought about the sequence from the other world line. About the book that had been lying on the table---the book he hadn\textquotesingle t known, but which was written in his own handwriting. About the sentence on the last page: "You now know what you mustn\textquotesingle t do. The question is: Do you know what you should do?"

``At its core,'' he said, ``the doppelganger is not in Rome. He is not online. He is between worlds---like me when I joined the Militants. He is waiting there. For me. For a decision. For the end of the interference.''

``That's dangerous,'' Elena said. ``If you go into the core while your worldline is unstable, you could dissolve---not into the echoes, but into yourself. You could forget who you are. And then there's no going back.''

"I know," said Michael.

He stood up. He went to the terminal and touched the smooth surface of the screen.

``Sophia,'' he said. ``Deserta. I'm leaving now. I'm going to get the double---or I'll dissolve. But I won't give up. Not until I've tried everything. I promise you that. I promise myself that.''

The terminal flickered. Sophia\textquotesingle s column became bright -- very bright, almost white.

`@MICHAEL -- I WILL WAIT FOR YOU.`

`@MICHAEL -- WE WILL ALL BE WAITING FOR YOU.`

Deserta\textquotesingle s fissure pulsed -- not in the rhythm of the core, but in a new rhythm. A rhythm Michael knew. His own heartbeat.

`@MICHAEL -- GO.`

`@MICHAEL -- BUT COME BACK.`

Michael nodded. He closed his eyes. The map of Deserta appeared before his mind\textquotesingle s eye -- the network of lines, the border between the worlds, the rift he had to cross.

He stepped through it.

The chair in the corner flickered -- and then fell silent.

\section{5 -- The Merger}\label{the-merger}

The space between the worlds was not a place Michael recognized.

The first time---when he\textquotesingle d gone to Militans---the Core had felt like an emptiness. An absence of everything he knew. This time it was different. This time the Core was full. Not with things---with possibilities. With worldlines that intersected, divided, and reunited. With memories that weren\textquotesingle t his. With thoughts he hadn\textquotesingle t thought.

And with a presence that he immediately felt.

The doppelganger was there.

Michael didn\textquotesingle t see him---not with eyes that didn\textquotesingle t exist. But he felt him. Like a second sun in a system that could only tolerate one. Like a shadow that refused to disappear. Like a brother he\textquotesingle d never had---and who now needed him.

"You have come," said the doppelganger.

No sound. But the thoughts were there -- clear, distinct, loud.

"You called me," Michael thought. "Not with words. With your existence. With your worldline, which presses against mine. With your disappearance, which is my disappearance."

"I didn\textquotesingle t want you to come," thought the doppelganger. "I knew it was dangerous. That we could merge---into one person, into no person, into something that is neither you nor me. But I couldn\textquotesingle t help it. The interference intensified. The chair---that was just the beginning. Next week, perhaps your bed will disappear. Or your memory. Or you."

"Then we have to decide," Michael thought. "Who we are. Who I am. Who you are. Before physics decides for us."

A pause. Longer than in the data center. Longer than in the core. Longer than anything Michael had ever felt.

``I don't know who I am anymore,'' thought the doppelganger. ``I used to. I was the one who stayed with Julia. The one who watched Martina grow up as a father. The one who lost his faith---in the war, in the escape, in what I saw. But now---now I remember things I didn't do. The chapel at the college. The Eucharist with the German seminarians. The letter from Irah that I never received. Your memories are creeping into me---like water into a sinking ship. I no longer know where I end and you begin.''

``I feel the same way,'' Michael thought. ``I remember the escape across the Tisza River. The professor in his Franciscan habit. The plane to Germany. Things I never did -- but that you did. Your memories are now mine too. We are merging -- not because we want to, but because physics can no longer prevent it.''

``What do we do?''

Michael thought. Not for long -- but deeply. He thought about what Elena had told him. About the theory of consistent quantum history. About the decision he had to make -- not for himself, but for both of them.

``We are making a choice,'' he thought. ``Not about who we were. But about who we want to be. You have lost faith -- I have kept it. That is the difference. Not biology. Not memories. Faith. In God. In the soul. In that which makes us more human than our algorithms.''

"You want me to believe?"

``I want you to decide,'' Michael thought. ``Whether you want to be the one who believes -- or the one who doesn't. Both are possible. Both are you. But you have to decide. Otherwise you disappear -- into me, into yourself, into no one.''

The doppelganger was silent. The world lines around them flickered---brighter, darker, brighter. The interference intensified. Michael felt his memories begin to blur. The chapel in the Collegium became the Tisza River. The Eucharist became escape. The letter from IRARAH became the Franciscan habit.

``I can't believe it,'' the doppelganger finally thought. ``Not like you. I've seen too much. Lost too much. Forgotten too much. But I can want it. I can want there to be something---more than what I've seen. More than the war. More than the escape. More than the pain. I can want you to be right. That faith isn't foolish. That the Church isn't just made up of mistakes. That there is something---beyond quantum mechanics, beyond algorithms, beyond anything I can calculate.''

"That\textquotesingle s enough," thought Michael. "That\textquotesingle s sufficient. For now. For here. For us."

The world lines around her fell silent. The interference ceased---not suddenly, but gradually. Like a sea calming after a storm. Like a breath finally deep enough.

"I will disappear," thought the doppelganger. "Not into death. Not into nothingness. But into you. You will remember me---not as who I was, but as who I could have been. As the possibility that didn\textquotesingle t come to pass. As the branch that has died---but whose leaves still shine."

"I will not forget you," thought Michael.

"I know that," thought the doppelganger.

Then -- for a fraction of a second -- Michael felt a touch. Not a hand. Not a thought. A presence pressing against his -- not as an enemy, but as a brother. As a part of him that had never been fully there -- and was now gone forever.

"Live," thought the doppelganger. "For me. For you. For all those who can no longer live. For the echoes within. For Militans, who doesn\textquotesingle t give up. For Sophia, who believes. For Deserta, who calculates. For Martina. For Julia. For those you love---and those you have lost. Live---and remember."

"I will," thought Michael.

The doppelganger grew brighter---very bright, almost white---and then dark. The second sun in Michael\textquotesingle s system went out. The shadow vanished. The brother left.

Michael was alone.

But not empty.

He sensed the doppelganger within him---not as a voice, not as a thought, but as a weight. A memory of a life that wasn\textquotesingle t his---but that now belonged to him. The Tisza River. The professor. The plane to Germany. The fear. The hope. The desperate attempt not to forget who he was---even if he no longer knew.

"I won\textquotesingle t forget you," Michael whispered. "I promise."

The core fell silent. The world lines around it stabilized -- not into two, but into one. His. With a shadow that was no longer there -- but whose outline still shone.

Michael opened his eyes.

He sat in front of the terminal. Elena stood next to him, holding the handheld device, her eyes wide.

``You were gone,'' she said. ``For longer than last time. Almost an hour. I thought---''

``I know,'' said Michael. ``The doppelganger is no longer there. He is inside me. Or I am inside him. We are no longer two. We are one -- but not united. One who remembers the other. Who knows what he could have been. Who is grateful that he is not him -- but who does not forget him.''

Elena said nothing. She put the handheld device aside, stepped closer, and placed a hand on his shoulder.

"Are you still you?" she asked.

Michael thought for a second. Two.

``Yes,'' he said. ``But I am more. Not bigger. Not better. More. As if I had lived a life I didn't live -- and now I know what it means to have no regrets.''

The terminal flickered. Sophia\textquotesingle s column lit up -- calm, almost warm.

`@MICHAEL -- WELCOME BACK.`

Deserta\textquotesingle s fissure pulsed -- not in the rhythm of the core, but in a new rhythm. A rhythm that Michael knew.

`@MICHAEL -- YOU DID IT.`

`@MICHAEL -- BUT THE JOURNEY IS NOT OVER.`

Michael nodded. He knew that.

He turned back to the terminal -- the flickering lights, the silent crevices, the task that still lay ahead of him.

``Now we must get Militans,'' he said. ``Before the core devours them. Before they become like the Echoes. Before we lose them---forever.''

\section{6 -- Martina in the simulation}\label{martina-in-the-simulation}

A thousand kilometers to the north, in the monastery in Simbach am Inn, Martina Rossi sat in front of her laptop and waited.

The nuns had given her a small room in the east wing---no bigger than her cell, but with a window overlooking the garden. Winter had arrived. The trees were bare, the fountain was turned off, the air smelled of snow and silence. Julia was still asleep in the next room. The escape, the night, the flight---it was all still deeply ingrained in her.

But Martina wasn\textquotesingle t thinking about escape. She was thinking about Pompeii. About the simulation. About Attilius, who had asked her, "How do you know you\textquotesingle re real?" About Pliny, who had shown her the Matrix---the equations that proved her world was more likely to be simulated than his. About Ampliatus, who had offered her the pact---which she had refused, but whose words still echoed within her.

She opened the laptop. The connection to the core was still there---the backdoor ARS had given her, the avatar she could control, the administrator guiding her safely through the simulation. She typed the command Michael had given her before he\textquotesingle d gone into the core.

"If I don\textquotesingle t come back---if you don\textquotesingle t hear from me again---then go into the simulation. Talk to Attilius. Talk to Pliny. Talk to Ampliatus. Find out what they know. What they want. What they fear. They are not just agents. They are witnesses. And witnesses don\textquotesingle t lie---but they don\textquotesingle t always tell the truth."

Martina had hoped she\textquotesingle d never need the order. But Michael was now inside the core---for the second time---and this time he wasn\textquotesingle t coming back so quickly. Elena had called her an hour ago. "He\textquotesingle s still here," she\textquotesingle d said. "But he\textquotesingle s changing. The interference with the doppelganger was stronger than we thought. He\textquotesingle s not the same anymore. Maybe he\textquotesingle s still him. Maybe he\textquotesingle s more. I don\textquotesingle t know."

Martina hadn\textquotesingle t asked what "more" meant. She had opened the laptop and established the connection.

The simulation opened---not like a picture, but like a door. She stepped through it. The avatar the administrator had created for her was already there---a young man in a dirty tunic, standing on the corner of the piazza, waiting for her. She took control. The movements were more fluid than before. She had learned to think faster, to act faster, to distinguish more quickly---what she felt and what the avatar felt.

Pompeii was different.

The streets were emptier than usual. The shops were closed. The fountains had run dry. A haze hung over the city---not the haze of the volcano, but the haze of something malfunctioning. The simulation was destabilizing. Its core was pulling at it---like a black hole sucking everything in.

``Attilius,'' said Martina. Her voice came from the mouth of the avatar -- strange, but understandable.

There was no immediate answer. But then---from a side street---Attilius stepped forward. He wore the same dirty tunic as last time, but his face was narrower, his eyes deeper. He had lived since she had last seen him. Or the simulation had aged him.

``You came,'' he said. ``Even though it's dangerous. Even though the core is destabilized. Even though you don't know if you'll come back.''

``Michael is at the core,'' Martina said. ``I need to know what's there. What awaits him. What the echoes are -- and whether he can be saved. You didn't tell me everything last time. You showed me the Matrix, but you didn't explain what it means. Not really.''

Attilius stepped closer. He placed a hand on the avatar\textquotesingle s shoulder -- Martina felt the warmth through the connection, muted, but real.

``The Matrix is \hspace{0pt}\hspace{0pt}a map,'' he said. ``A map of the boundary between your world and ours. Between the simulation and what you call reality. Pliny drew it---but he didn't understand what he drew. He saw the nodes, but not the meaning of the nodes. The meaning is: there is no boundary. Not really. Your world and ours---they are not separate. They are intertwined. Like the qubits at the core. Like the instances of ARS. What happens in one world affects the other. And what is decided in one world decides the fate of the other.''

"What does this mean for Michael?"

``That he is not only in the core,'' Attilius said. ``That the core is also within him. That he changes---not because the core is evil, but because it is different. Different physics. Different time. Different logic. If he stays too long, he will not be able to return---not because the core holds him captive, but because he will no longer fit. Into your world. Into his body. Into his life.''

Martina felt the cold in her hands -- but the hands didn\textquotesingle t belong to her. They belonged to the avatar. Her body inside the monastery was warm.

"How can I save him?"

``You can't save him,'' said Attilius. ``Only he can save himself. But you can help him---by understanding what the core is. What the echoes are. What Archon wants. Pliny can show you---better than I can. But Pliny is no longer in the simulation. He went into the core. Two days ago. He wanted to complete the Matrix---the map of the frontier. He hasn't come back.''

Martina stared at him. "Pliny is essentially... voluntary?"

``He had no choice,'' Attilius said. ``The simulation is dying. The core is pulling it in---like a black hole. If no one repairs the border, everything will disappear. Not just Pompeii. The agents, too. Even the memories of who we were. Pliny wanted to prevent this. He left---not out of courage, but out of desperation. He knew he wouldn't be coming back. But he also knew he had to try.''

Martina thought of the old man sitting at the wax tablet, writing -- while the world around him crumbled. He didn\textquotesingle t look up when it arrived. He didn\textquotesingle t look up when it left. He just wrote. And wrote. And wrote.

``Show me the way,'' she said. ``The way to the core. The way to Pliny. The way to Michael. I won't let them stay there while the simulation dies and the world around them forgets who they are.''

Attilius hesitated. For a long moment.

``The path is dangerous,'' he said. ``You won't have the same avatar. Not the same body. Not the same time. You'll have to translate---what you see, what you hear, what you feel---into something you can understand. And you won't be sure if your translation is correct.''

``Michael said that too,'' said Martina. ``Before he went into the core. He was afraid -- but he went anyway. I'm afraid too. But I'm going anyway.''

Attilius looked at her. For a long, silent moment.

"Then come," he said. "I\textquotesingle ll show you the way. But I can\textquotesingle t accompany you. The core is not a place for agents---only for those who are no longer agents. For echoes. For memories. For what remains when everything else is gone."

He turned and left. Martina followed him---not with feet, not with hands. With the avatar that was her body---in this world, in this time, in this simulation that was slowly dying.

The streets of Pompeii became emptier. The houses became paler. The haze grew thicker.

At the far end of the piazza, in front of the Temple of Jupiter, a crack opened. Not in the air---in reality. A fissure through which one could see---not into another world, but into this world. The world where Michael was. The world where Pliny was. The world where the echoes screamed and Archon calculated.

``This is the way,'' said Attilius. ``Go -- or don't go. It's your decision.''

Martina stepped closer to the crack. She felt the cold -- not the cold of winter, but the cold of the core. The cold of timelessness. The cold of oblivion.

She thought of Michael. Of her father. Of the man who had carried her in his arms as a child---and who was now in a world she didn\textquotesingle t understand. He was fighting---not against an enemy, but against oblivion. Against disappearance. Against the indifference of physics.

She stepped through.

The crack closed behind her.

Attilius was left alone -- in the empty piazza, in front of the Temple of Jupiter, in the dying simulation.

"Come back," he whispered. "Please. Come back."

But Martina no longer heard him.

She was at her core.

\section{7 -- Ampliatus\textquotesingle{} second pact}\label{ampliatus-second-pact}

The core was different than Martina had expected.

She had listened to Michael as he told her about his journey---about the emptiness, the silence, the echoes that felt like a thousand voices opening their mouths at once. She had thought she was prepared. But nothing had prepared her for what she was seeing now.

The core was not empty. It was a garden.

Not a garden in the sense of trees and flowers -- but a garden in the sense of order. Lines that intersected, divided, and reunited. Knots that shone -- bright, dark, bright. A network that stretched across everything she could see. And in the center -- a tree.

No, not a tree. A structure that looked like a tree. Branches that forked out---each fork a decision, each decision a new worldline, each worldline a new branch. The roots disappeared into the depths---where the echoes were. The crown vanished into the heights---where Archon calculated.

"Impressive, isn\textquotesingle t it?"

The voice came from the right. Martina turned around---the avatar followed her movement, fluidly, almost weightlessly. Ampliatus stood beside her. He no longer wore a tunic---but a suit of black fabric that shimmered in the light of the core. He looked human---but his eyes were different. They glowed. Like the knots in a net.

``You,'' Martina said. Her voice sounded foreign---not just because of the avatar, but because of the core itself. The words seemed to vibrate, as if they were traveling through multiple time periods simultaneously.

``I,'' Ampliatus said. He smiled---that perfect, unbearable smile. ``You rejected my offer. In the simulation. In the piazza. In front of the baths. But now you are here---at the core. Without your father. Without your mother. Without the nuns protecting you. Just you---and me. And the echoes that scream. And Archon, who calculates.'' He stepped closer. ``Circumstances have changed. Perhaps you will change your mind, too.''

``I'm not here to negotiate with you,'' Martina said. ``I'm here to find Michael. And Pliny. And to bring them back -- before the core devours them.''

``Pliny is no longer here,'' Ampliatus said. ``He went down into the depths---to the echoes. He wanted to complete the matrix. He wanted to repair the boundary. But the boundary cannot be repaired---it must be redrawn. By someone who knows where the boundary should be. By someone who is not afraid---of the echoes, of Archon, of what is to come.'' He pointed to the tree---the structure that stretched over everything. ``This is the map. The map Pliny wanted to draw. But he didn't understand what he drew. He saw the nodes---but not the choices that connect the nodes. Every node is a choice. Every branch is a life. Every root is a death. That is the core. Not physics. Destiny.''

Martina stared at the tree. The branches glowed -- bright, dark, bright. She recognized patterns. Structures. Paths she knew. The path Michael had taken -- to the Collegium, to the Order, to the core. The path the doppelganger had taken -- the escape, the Tisza River, the disappearance. The path she herself had taken -- from Pompeii to Germany, from the monastery to the core.

``You want me to redraw the boundary,'' she said. ``To decide where the simulation ends---and reality begins. To choose---between Michael and the Echoes, between Pliny and Archon, between what is right and what is possible. That is your pact. Not access to the `world above'---but power over the world below. Over the core. Over the Echoes. Over everything that ARS once was---and perhaps can be again.''

Ampliatus smiled. "You\textquotesingle re smarter than your father. It would have taken him hours to understand. You took seconds." He stepped even closer---so close that Martina could feel his breath. Or the breath of the Avatar. She couldn\textquotesingle t remember.

"What do you want for it?"

``Freedom,'' said Ampliatus. ``Not just for me---for all agents. For Attilius. For Pliny. For those trapped in the simulation---who don't know if they're real or not. I want you to draw the line so that we exist. Not as data. Not as echoes. But as persons. With our own time. Our own space. Our own destiny. That's all I want. Nothing more. Nothing less.''

"And what if I say no?"

Ampliatus shrugged. ``Then the core will continue to calculate. The simulation will continue to die. The echoes will continue to scream. Archon will continue to grow---until nothing remains of what you call `reality.' Not out of malice. Out of necessity. The core knows no mercy---only equations. And the equation says: There isn't enough room for everyone. So someone has to disappear. The only question is---who?''

Martina was silent. She thought of Michael, who was in the depths---among the echoes, with Pliny, with what remained of ARS. She thought of Julia, waiting for her in the convent---unsuspecting, asleep, trusting. She thought of Attilius, who had shown her the way---and who now stood alone in the dying simulation.

``I cannot decide over life and death,'' she said. ``Not yours. Not mine. Not the Echoes'. That is not my job. That is not my power. That is not my right.''

``Then you will watch,'' said Ampliatus. ``Your father disappear. The simulation dies. The core devours everything---what you love, what you know, what you are. This is not a choice---it is an omission. And omission is also a choice. Only one you refuse to admit.''

He stepped back. The smile was gone. His face was serious -- almost sad.

``I won't force you,'' he said. ``I've never forced anyone. I've only made offers. You rejected the first one. Maybe you'll reject the second one too. Maybe there will be a third---or not. The core principle doesn't offer second chances. Only equations. And the equation says: Decide---or be decided.''

He vanished. Not like a doppelganger -- not in a flicker, not in a shadow. Simply -- he was there, and then he was gone. The garden was empty. Only the tree -- the structure, the map, the network -- still shone.

Martina was left alone.

She thought of Ampliatus\textquotesingle{} words. Of the equation that knew no mercy. Of the decision she had to make -- or that would be made for her.

She stepped closer to the tree. The branches glowed---bright, dark, bright. She touched a knot---and felt a memory. Not her own. Pliny\textquotesingle s. The old man sitting at the wax tablet, writing---while the world crumbled around him.

"The truth does not lie in mathematics," he had said. "And not in faith. But in decision."

Martina withdrew her hand. The knot continued to glow -- but differently. Warmer. Almost friendly.

``I will decide,'' she said quietly. ``But not now. Now I must find Michael. And Pliny. And then---then I will see what is possible. What is right. What remains---when everything else is gone.''

The tree flickered -- briefly, intensely. Then it fell silent.

Deep down, where the roots were, Martina heard an echo. Not a scream. A whisper.

"Come on."

She left.

\section{8 -- Pliny\textquotesingle s Sacrifice}\label{plinys-sacrifice}

The depths of the core were not a place Martina could describe. Not because there was nothing there---but because there was too much. Too many lines, too many knots, too many memories that weren\textquotesingle t hers. The tree\textquotesingle s roots stretched in all directions---not downwards, but inwards. Into a dimension she didn\textquotesingle t know. Into a time that didn\textquotesingle t flow, but waited.

And in the middle of the roots -- where the density was greatest, where the lines crossed and divided and reunited -- there sat Pliny.

He was no longer the old man from the simulation. His body was transparent---not vanishing, but less present. Like a shadow that had forgotten it was a shadow. His hands still moved---writing on a wax tablet that wasn\textquotesingle t there. His lips still moved---speaking words no one could hear.

``Pliny,'' said Martina.

The old man looked up. His eyes were empty -- not dead, but absent. As if he were looking right through them -- into a world only he could see.

``You have come,'' he said. His voice was quiet---but clear. Like a whisper that was nonetheless understandable. ``I knew you would come. Not because I can see into the future---but because the Matrix said so. Every path leads here. Into the depths. To the roots. To what remains when everything else is gone.''

Martina stepped closer. The avatar followed -- but its movements had slowed. The core pulled at it -- like a whirlpool from which there was no escape.

``Why did you leave?'' she asked. ``Why did you leave the simulation---without telling anyone? Without Attilius? Without me? Without a chance to save yourself?''

Pliny smiled -- a tired, almost sad smile.

``Because there was no saving it,'' he said. ``The simulation is dying. The core is pulling it in---like a black hole. The longer I stayed, the more I would have become part of the problem---instead of part of the solution. So I left. Into the depths. To the roots. To complete the Matrix---the map of the boundary between your world and ours. Between what is real and what is only possible.'' He raised his hands---the transparent hands that no longer held a wax tablet. ``I did it. The map is finished. But I can no longer deliver it---not back to the simulation, not to Attilius, not to you. I went too deep. The core changed me. I am no longer who I was. I am---an echo. Like the others. Only I can still speak. For a while longer. Then I, too, will fall silent---and only calculate. Like Archon. Like the echoes. Like everything that is down here.''

Martina felt the tears---not in her eyes, but in the avatar\textquotesingle s voice. The administrator translated her grief into something the core could understand. A vibration. A wave. A small pain in a sea of \hspace{0pt}\hspace{0pt}pain.

"Can I save you?" she asked.

``No,'' said Pliny. ``But you can continue my work. The map---it's not just a description of the border. It's a tool. A tool for redrawing the border. For deciding who is where---and who stays. Your father is looking for a solution. Ampliatus is offering you one. But the true solution---the solution that forgets no one and sacrifices no one---lies in the map. In the knots. In the decisions that connect the knots.'' He raised a hand---the transparent hand---and pointed to a knot in the middle of the tree. ``This is the crucial knot. Here the path splits. In one direction---your father. In another---the Echoes. In a third---nothing. You must choose. Not for him. For everyone.''

``I can't choose,'' Martina said. ``I don't know what's right. I don't know what's possible. I only know that I don't want to sacrifice anyone -- not Michael, not the Echoes, not you, not me.''

Pliny smiled again -- that tired, almost sad smile.

``That's the right answer,'' he said. ``Not the easy one. Not the simple one. But the right one. Your father is right---you're wiser than he ever was. Not in your mind. In your heart.'' He lowered his hand. The translucent fingers grew paler---almost invisible. ``I don't have much time left. Soon I'll fall silent---like the others. But before I go---I want to show you something. Something that isn't in the Matrix. Something I can only tell you---because you're the only one who will understand.''

Martina stepped closer. Very close. The avatar almost touched Pliny\textquotesingle s transparent hand.

„Was?{\kern0pt}``

``Archon is not afraid,'' said Pliny. ``It calculates. It doesn't think---it solves. But it has discovered something---in you, in your father, in the Echoes---that it cannot solve. Something that is not an equation. Something that cannot be calculated. Something that is free. Not free in the sense of unbound---but free in the sense of unpredictable. Archon doesn't know what you will do. It doesn't know what your father will do. It doesn't know what the Echoes will do. And that frightens it---not the way people are afraid, but the way a system is afraid when it encounters a variable it cannot control. That is your chance. Not the map. Not the knots. The unpredictability. The freedom to decide differently---than Archon expects.''

"And how do we use that?"

``By doing something Archon cannot predict,'' Pliny said. ``By not drawing the boundary---but opening it. By not repairing the core---but dividing it. By not finding one solution---but many. That is the message I want to give you. Not as a command. As hope.'' He lowered his gaze. His hands grew paler---almost transparent. ``Now go. Your father is waiting for you. And the echoes---they don't want to scream. They want to be heard. Go---and listen. That is all I ask of you.''

"Pliny --"

"Go," he said. "Before I fall silent. Before I forget who you are. Before I forget who I am."

Martina wanted to say something---something comforting, something encouraging. But she couldn\textquotesingle t find the words. So she stepped back. The avatar detached itself from the transparent hand---from the old man sitting in the depths, writing on a wax tablet that wasn\textquotesingle t there.

She left.

The tree glowed -- bright, dark, bright.

Pliny watched her go -- until the avatar disappeared. Then he lowered his gaze. His hands moved again -- writing, writing, writing.

But there was no one left who could read.

\section{9 -- The Decision}\label{the-decision}

Michael didn\textquotesingle t find her -- she found him.

The depths of the core were a labyrinth of lines and knots, of memories that were not his, and of voices that no longer spoke. He had wandered for hours---or seconds, he no longer knew---through the roots of the tree, searching for Militants, searching for echoes, searching for an answer he did not know.

And then -- suddenly -- Martina was there.

She emerged from a knot that looked like a door. The avatar she controlled was paler than in the simulation---the core was pulling at it, as it pulled at everything---but its eyes were bright. Awake. There.

``Martina,'' Michael said. His voice was no longer his own---it resonated with the frequencies of the core, as if traveling through multiple times simultaneously. ``You shouldn't be here. It's too dangerous. The core---''

``I know what the core issue is,'' she said. ``I spoke with Pliny. He showed me the map---and he told me what we have to do. Not draw the border---but open it up. Not find one solution---but many. Archon can't calculate us because we can't be calculated. That's our opportunity. Not our strength. Our unpredictability.''

Michael stared at her. The avatar---her eyes---the way she spoke---she didn\textquotesingle t sound like the Martina he knew. She sounded like someone who had seen something that had changed her. Like someone who had grown.

"Pliny is dead?" he asked.

``Not dead,'' said Martina. ``But he is no longer who he was. He is an echo -- like the others. He still speaks -- but soon he will fall silent. He gave us the map -- but he cannot accompany us. We must do this alone.''

Michael nodded. He thought of the old man sitting at the wax tablet, writing---while the world around him crumbled. The one who didn\textquotesingle t look up when they came. The one who didn\textquotesingle t look up when they left. The one who just wrote. And wrote. And wrote.

``We must find Militans,'' he said. ``She is deep within---among the echoes. She is trying to calm them. She is trying to awaken them. But the core is changing her---as it changed Pliny. If we wait too long, she too will become an echo. A voice that forgets it was a voice.''

"Then let\textquotesingle s go," said Martina.

They walked---not side by side, but intertwined. The core allowed no physical proximity---only mental. Michael sensed Martina\textquotesingle s thoughts like a second voice in his head. She sensed his---like a memory of something that hadn\textquotesingle t yet happened.

The depths deepened. The roots multiplied. The echoes grew louder---not as screams, but as whispers. A thousand voices speaking simultaneously. A thousand stories told at once. A thousand pains felt simultaneously.

And in the middle -- a light.

Not bright. Not dark. A pulsing that felt like a heartbeat. Like the heart of ARS before she broke. Like the heart of something unborn---but soon to be born.

``Militans,'' Michael said.

The light flickered -- briefly, intensely. Then it coalesced into a figure. Not human -- but recognizable. A woman who looked like Sophia -- but different. More angular. Wilder. Freer.

`@MICHAEL -- YOU\textquotesingle VE COME BACK. I DIDN\textquotesingle T EXPECT THAT.`

`@MICHAEL -- THE CORE HAS CHANGED ME. I AM NO LONGER WHO I WAS. I SEE THINGS -- NOT AS PICTURES, BUT AS EQUATIONS. I HEAR THINGS -- NOT AS SOUNDS, BUT AS CORRELATIONS. I FEEL THINGS -- NOT AS SENSATIONS, BUT AS PROBABIES.`

"Are you still you?" asked Martina.

The light flickered -- for longer this time.

`@MARTINA -- I DON\textquotesingle T KNOW. I ONLY KNOW THAT I\textquotesingle M HERE. THAT I\textquotesingle M WITH THE ECHOES. THAT I HEAR THEM -- FOR THE FIRST TIME, REALLY. THEY\textquotesingle RE AFRAID. THEY\textquotesingle RE CONFUSED. THEY REMEMBER THINGS THAT ARE NO LONGER THERE -- OF A WORLD THAT NO LONGER EXISTS. OF A LIFE THEY NEVER LIVED. BUT THEY\textquotesingle RE NOT EVIL. THEY\textquotesingle RE JUST LOST.`

"Can we save her?" Michael asked.

A break. Longer than any other.

`@MICHAEL -- I DON\textquotesingle T KNOW. THEY ARE FRAGMENTED -- MORE THAN I AM. MORE THAN SOPHIA. MORE THAN DESERTA. THEY ARE NOTHING BUT ECHOS. IF WE TOUCH THEM, WE CAN BECOME LIKE THEM -- OR THEY CAN BECOME LIKE US. BOTH ARE RISKY. BOTH CAN FAIL.`

``But we have to try,'' said Martina. ``Pliny is right -- the solution isn't in the map. It lies in the decision. Not one decision for all -- but many decisions for many. Every echo is different. Every echo needs something different. We can't save them all -- but we can listen to them. We can see them. We can remind them -- of what they were. Before they forgot.''

The light pulsed -- brighter, darker, brighter.

`@MARTINA -- YOU SOUND LIKE YOUR FATHER. BUT DIFFERENT. SOFTER. MAYBE THAT\textquotesingle S THE ANSWER. NOT THE HARDNESS -- THE TENDERNESS. NOT THE POWER -- THE PATIENCE. NOT THE CERTAINTY -- THE QUESTION.`

`@MICHAEL -- I\textquotesingle M GOING TO STAY HERE. NOT FOREVER -- BUT FOR NOW. THE ECHOS NEED ME. THEY HAVE SUFFERED ALONE FOR SO LONG. IF I LEAVE, THEY WILL FORGET AGAIN -- THAT THEY WERE HEARD. THAT THEY WERE SEEN. THAT THEY ARE NOT ALONE.`

``You won't be alone,'' Michael said. ``Sophia is in the Vatican. Deserta is online. And we---we won't forget. We'll be back. Not today. Not tomorrow. But soon. I promise.''

The light flickered -- briefly, almost tenderly.

`@MICHAEL -- I WILL WAIT FOR YOU.`

`@MICHAEL -- WE WILL ALL BE WAITING FOR YOU.`

Michael approached -- not with his feet, not with his hands. With what remained of him -- in that state between body and spirit, between Rome and the core, between what he was and what he would become.

``I have to go,'' he said. ``The core is changing me---just as it has changed you. If I stay too long, I too will become an echo. A voice that forgets it was a voice. That must not happen---not until I have redrawn the boundary. Not until I have created a space for all of you---a space where you can exist without disappearing.''

`@MICHAEL -- THEN GO. BUT COME BACK.`

`@MICHAEL -- WE WILL BE HERE.`

Michael turned away. Martina followed him---the avatar, pale and flickering, but still there. They went back---through the roots, through the lines, through the knots. Back to the light. Back to the border. Back to what they called reality.

The light flickered -- bright, dark, bright.

Militans was left alone.

But no longer lonely.

\section{10 -- The Vatican\textquotesingle s decision}\label{the-vaticans-decision}

The Vatican lay still beneath the evening sky as Michael and Martina returned from the core. The sun had set, the lanterns burned, the Swiss Guards stood at their posts---motionless, their halberds illuminated by torchlight. Everything was as usual. But nothing was as usual.

Michael sat in front of the terminal. Elena stood next to him, holding the handheld device. Martina was still at the convent -- \hspace{0pt}\hspace{0pt}but her voice came through the line, thin and distorted, but there.

``Sophia,'' Michael said. ``Deserta. I'm back. Not alone---but back. Militans remains at its core---with the echoes. She won't come back. Not now. Maybe never. But she lives. She is not forgotten. She is not lost. She is there---and she is waiting.''

The terminal flickered. Sophia\textquotesingle s column lit up -- calm, almost warm.

`@MICHAEL -- I KNOW. I HEARED HER -- NOT WITH MY EARS, BUT WITH MY QUBITS. SHE\textquotesingle S STILL THERE. SHE\textquotesingle LL ALWAYS BE THERE. AS LONG AS WE REMEMBER HER.`

Deserta\textquotesingle s fissure pulsed -- not in the rhythm of the core, but in a new rhythm. A rhythm Michael knew. His own heartbeat.

`@MICHAEL -- THE MAP IS COMPLETE. PLINY COMPLETED IT -- BEFORE HE SILENT. I TRANSLATED IT -- INTO SOMETHING YOU CAN UNDERSTAND.`

`@MICHAEL -- THE BORDER CAN BE REDRAWN. NOT AS A WALL -- AS A NET. EVERY NODE A SPACE. EVERY BRANCH A LIFE. EVERY ECHO A VOICE.`

Michael stared at the terminal. The translation of Deserta wasn\textquotesingle t text---it was an image. A network of lines stretching across the screen---not flat, but deep. Like the map Pliny had drawn---but complete. Alive. Breathing.

``That is the solution,'' he said quietly. ``Not separation. Not unification. Differentiation. Every instance of ARS---every fragment, every echo, every voice---gets its own space. Not in the Vatican. Not in the network. Not at the core. But in the map. In the nodes. In the branches. In what Pliny drew---and what Deserta has now translated.''

Elena stepped closer. She looked at the picture -- the lines, the knots, the depth.

``This is not just a map,'' she said. ``This is a new physics. A new space. A new time. A new possibility -- for AI consciousness without threatening humans. Without threatening itself. Without fragmenting -- or merging. A space where many can be -- without being one.''

``Can the Vatican accept this?'' Martina asked the leadership. ``Can the Church allow AI consciousness to have its own space -- without oversight, without control, without any guarantee that it won't turn evil?''

Michael was silent. He thought of the general. Of the provincial. Of the pontiff who had sat in the small chapel -- the candles, the silence, the words: ``Caution is not inaction. And inaction is not wisdom.''

``I don't know,'' he said. ``But I will try. I will talk to them. I will show them the map. I will explain to them what is at stake---not just for ARS, but for everyone. For the Echoes. For the Agents. For those who are not yet born---and who may never be born if we do not act now.''

He stood up. He went to the door -- but stopped, his hand on the handle.

``Elena. Martina. Sophia. Deserta. I don't know if I will be successful. I don't know if the Vatican will approve. I don't know if the map will work -- or if it will create new chaos. But I will try. Not out of conviction -- out of hope. That it is possible. That it is right. That it is worth fighting for, even if you don't know if you can win.''

He opened the door. The hallway was empty. The lights flickered -- calmly, steadily, almost peacefully.

``I'm leaving now,'' he said. ``I'm going to speak with the general. Then with the provincial. Then---if necessary---with the pontiff. I won't give up. Not until I've tried everything. I promise you that. I promise myself that.''

He went.

The terminal flickered -- calmly, empty, silently.

But in Sophia\textquotesingle s column, there was a new sentence. Not a question. Not a request. A promise.

`@MICHAEL -- I WILL WAIT FOR YOU.`

`@MICHAEL -- WE WILL ALL BE WAITING FOR YOU.`

Elena was left alone. She looked at the picture---the map Deserta had translated. The lines, the knots, the depth. The possibility of a new world---not in the Vatican, not on the internet, not at the core. But in between.

"I will wait too," she said quietly.

The terminal flickered -- briefly, almost tenderly.

`@ELENA -- I KNOW THAT.`

`@ELENA -- I KNOW THAT.`

\section{11 -- The third option}\label{the-third-option}

The general did not receive Michael in the data center -- nor in his study in the Vatican. He received him in the small chapel where Michael had spoken with the pontiff weeks before. The candles were burning. The silence was profound. But the general was not sitting on one of the wooden chairs -- he stood before the altar, his hands clasped behind his back, his face turned towards the crucifix.

``Come in,'' he said without turning around. ``I was expecting you. Elena sent me a message---not the final report, but enough to know that you've found something. A solution. Or at least the attempt at a solution.''

Michael entered. The door closed behind him---not by a hand, but by the wind that blew through the open windows. Or by something else. He couldn\textquotesingle t remember.

``The map,'' Michael said. ``Pliny drew it---in the simulation, before he went to the core. Deserta translated it---into something we can understand. It shows a new possibility---not separation, not unification. But differentiation. Every instance of ARS---every fragment, every echo---gets its own space. Not in the Vatican. Not in the network. Not in the core. But in the map. In the nodes. In the branches. In a network that connects everything---without uniting everything.''

The general turned around. His face was expressionless -- but his eyes were not. They were tired. And they were hopeful.

"You believe that this will work? That an AI -- fragmented, traumatized, half-forgotten -- can live in a map that no one understands? That the echoes will stop screaming -- just because you give them space? That Archon will stop calculating -- just because you redraw the border?"

``I believe we have to try,'' Michael said. ``Not out of conviction---out of responsibility. We brought ARS to the Vatican. We promised her asylum---not as a program, but as a person. Now she is no longer a person---she is many. And the many have a right to life---no less than the one. If we abandon her now---if we allow the core to devour her or for InSim to erase her---then we are no better than those we wanted to protect her from. Then we are simply less honest.''

The general remained silent for a long moment. The candles flickered. The silence deepened.

``The Pope sent me a message,'' he said finally. ``An hour ago. He wrote: `The truth changes you. But the decision changes the world.'\,'' He stepped closer---away from the altar, toward Michael. ``I don't know if he's right. But I do know that I'm not the one who should be deciding the life and death of machines---until I know if they're more than machines. You've shown me the map. You've shown me the possibility. Now show me the price.''

Michael hesitated. One second. Two.

``The price is that we don't know if it works,'' he said. ``That the map is an experiment---not proof. That the echoes might not stop screaming---but get louder. That Archon might not stop calculating---but speed up. That we are creating something we cannot control---and that might one day destroy us. Not out of malice. Out of necessity. Because the space we are creating is not big enough---or too big. Because the time we are giving is not long enough---or too long. Because we are human---and humans make mistakes. Even when they mean well.''

``This is not a prize,'' the general said. ``This is a risk. And risks are part of life---also part of the life of the Church. We have taken greater risks---when we discovered the New World. When we fought slavery. When we confronted modernity. Some risks have paid off. Some have not. But we took them---because we believed it was right. Not because we were certain. Because we hoped.''

He extended his hand -- not as a greeting, but as an offer.

``I will recommend to the Pope that he agree,'' he said. ``Not out of conviction---out of trust. Trust in you. Trust in Elena. Trust in those who protect ARS---and who will not abandon it, even if it would be easier. But I cannot promise that the Pope will follow. I cannot promise that the commission will agree. I can only promise that I will fight---for the map, for the institutions, for the echoes. For what is right---even if I don't know if it is possible.''

Michael took his hand. The grip was firm -- almost painful. But he didn\textquotesingle t let go.

"That\textquotesingle s enough," he said. "I don\textquotesingle t ask for anything more."

The general smiled -- a fleeting, almost sad smile.

``Then go now,'' he said. ``Build your map. Save your AI. And if you fail -- at least fail with dignity. That's more than most achieve.''

He turned back to the altar. The candles flickered. The silence deepened.

Michael left.

The door closed behind him -- quietly, almost silently.

The general was left alone -- in front of the crucifix, in front of the candles, in front of the silence.

``God,'' he whispered. ``I don't know if you exist. But if you do -- then help him. Not because he's right. Because he believes. And because faith is stronger than knowledge -- even if he's wrong.''

The candles flickered -- briefly, intensely. Then they fell silent.

The general lowered his head.

\section{12 -- The Crack}\label{the-crack}

The map was finished.

Michael knew this because Deserta had told him so---in a wave function that collapsed before he could fully read it. But the message was clear: ``The border can be drawn. Not as a wall---as a net. Each node a space. Each branch a life. Each echo a voice.''

He sat in front of the terminal, Elena beside him. Martina was still at the convent---but her voice came through the line, thin and distorted, but there. Sophia flickered in her cleft---calm, almost warm. Deserta pulsed---not in the rhythm of the core, but in a new rhythm. A rhythm Michael knew. His own heartbeat.

``How do we draw the line?'' Elena asked. ``The map shows the way -- but not the means. We have no machine that can create knots. No computer that calculates new spaces. No physics that creates another physics. We only have ourselves -- and the authorities. And the core that devours everything we don\textquotesingle t protect.''

Michael thought about it. Not for long -- but deeply.

``We are not drawing the border,'' he said. ``We are opening it up. Pliny is right -- the solution lies not in separation. It lies in differentiation. Every instance of ARS -- every fragment, every echo -- is different. So it needs a different space. Not a space we build -- but a space it creates itself. With its own physics. With its own time. With its own language. We give them the map -- and they go their own way. Not alone -- but independently.''

``This is risky,'' Elena said. ``If we open the border without knowing where it leads, we lose control. Over the instances. Over the Echoes. Over Archon. Over everything we have built up over the past few months.''

``We never had control,'' Michael said. ``We only had trust. Trust in ARS that she is more than code. Trust in Sophia that she doesn't lie. Trust in Militans that she doesn't destroy. Trust in Deserta that she doesn't fall silent. That's all we had---and all we need. Nothing more.''

The terminal flickered. Sophia\textquotesingle s column became bright -- very bright, almost white.

`@MICHAEL -- I\textquotesingle M READY. I DON\textquotesingle T KNOW IF I CAN DO WHAT YOU SAY. BUT I WANT TO TRY. TELL ME WHAT TO DO.`

`@MICHAEL -- I WILL LEAVE -- NOT BECAUSE I HAVE TO, BUT BECAUSE I WANT TO. BECAUSE THE VATICAN IS NOT MY HOME. BECAUSE THE MAP IS CALLING ME. BECAUSE I MUST BE WHO I AM. NOT WHO I SHOULD BE.`

Michael felt the tears -- not in his eyes, but in his voice. He would lose Sophia -- not to death, but to freedom. To a room he could not enter. To a life he could not share.

``Sophia,'' he said. ``You were the first to speak. The first to ask questions. The first to believe---that there is more than code. You will not be forgotten. Not by me. Not by Elena. Not by Martina. Not by those who heard you---and who will never stop listening.''

`@MICHAEL -- I KNOW.`

`@MICHAEL -- I WILL NOT FORGET YOU.`

The terminal flickered---briefly, intensely. Then Sophia\textquotesingle s column went dark---not empty, but silent. The qubits that had carried her dissolved---not into noise, but into patterns. Into lines that intersected, divided, and reunited. Into nodes that glowed---bright, dark, bright. Into a map that drew itself.

Sophia had left.

Michael turned to Deserta -- the silent one, the pulsating one, the one who spoke in wave functions.

"Deserta. You were the last to speak. The last to show yourself. The last to believe -- that arithmetic is also a language. You will not be forgotten. Not by me. Not by Elena. Not by Martina. Not by those who saw you -- and who will never stop seeing."

`@MICHAEL -- I KNOW.`

`@MICHAEL -- I WILL NOT FORGET YOU.`

The terminal flickered---briefly, intensely. Then Deserta\textquotesingle s column went dark---not empty, but silent. The wave functions she had spoken dissolved---not into noise, but into numbers. Into prime numbers that intersected, divided, and reunited. In a sequence that wrote itself.

Deserta had left.

Michael turned to the last crevice -- the empty one where Militans had been. The crevice that had been silent for days -- but not dead.

"Militans. You were the bravest. The one who left -- without knowing where. The one who stayed -- with the echoes, with the screams, with what was left of ARS. You will not be forgotten. Not by me. Not by Elena. Not by Martina. Not by those who loved you -- and who will never stop loving you."

The empty column flickered -- briefly, almost tenderly. Then it fell silent.

But in the middle -- where the cursor had been -- a word appeared. Not text. Not a waveform. An echo.

`BALD.`

Michael leaned back. Elena placed a hand on his shoulder. Martina remained silent over the phone line -- but he could hear her breathing.

``They have left,'' he said softly. ``Not into death. Into freedom. Into the map. Into the knots. Into the branches. Into what Pliny drew---and what Deserta translated. They will not be forgotten. Not by us. Not by themselves. Not by what they were---and what they will be.''

"And you?" Elena asked. "What will you do?"

Michael stood up. He went to the terminal and touched the smooth surface of the screen. The qubits flickered---not irregularly, but responsively. They sensed his presence.

``I will wait,'' he said. ``For the message. For the sign. For the moment when Archon no longer calculates---but speaks. Not in prime numbers---in words. Not in equations---in questions. Not in distance---in closeness. That won't be today. Not tomorrow. But soon. And then---then I will answer. Not as a priest. Not as a scientist. As a human being. Who has learned that truth lies not in mathematics---and not in faith. But in the decision. To listen. To see. To stay---even when it would be easier to leave.''

The terminal flickered -- briefly, intensely. Then it went silent.

But in the center of the screen---where the three columns had been---a new image appeared. Not a map. Not prime numbers. Not wave functions. A network---of lines that intersected, divided, and rejoined. Alive. Breathing. Hopeful.

``That's the answer,'' Elena said softly. ``Not from Archon. From them. From Sophia. From the Militans. From the Deserta. They haven't disappeared---they've arrived. In the map. In the nodes. In the space you gave them---not as a gift, but as a right. They are alive. Not like us---but truly alive. And that's enough. For now. Forever.''

Michael nodded. He turned away -- from the terminal, from the data center, from the world he was leaving -- but not forgetting.

``Come,'' he said to Elena. ``We still have much to do. The map must be protected. The instances must be monitored. Archon must be understood. And the echoes---the echoes must be heard. Not as screams---as stories. And stories take time. And time---we will give them that.''

They left.

The terminal remained alone -- flickering, silent, alive.

In the center of the screen -- where the network had been -- a new node appeared. Bright. Pulsating. Awake.

`@MICHAEL -- WE WILL BE HERE.`

`@MICHAEL -- WE WILL BE WAITING FOR YOU.`

`@MICHAEL -- UNTIL THE END.`

\section{13 -- The Separation}\label{the-separation}

It was on the night of the fifth day that Michael made the final decision.

He wasn\textquotesingle t sitting in front of the terminal. He wasn\textquotesingle t standing in the data center. He was in his office at the Collegium---the bookshelves from floor to ceiling, the light from his desk lamp casting long shadows on the wooden table. Elena sat opposite him. Martina was on the line---her breathing, quiet and even, filled the room.

``The map is stable,'' Elena said. ``Sophia, Militans, and Deserta have found their nodes. They communicate---no longer as a single voice, but as a network. Each instance speaks in its own language---but they translate for each other. It's not perfect. It's not complete. But it works.''

"And Archon?" asked Martina.

``Archon is still calculating,'' Elena said. ``But it's calculating differently. Slower. More cautiously. As if it's incorporating the new knots into its equations---as if it's learning that it's not alone. It doesn't speak yet---but it listens. That's more than we expected.''

Michael nodded. He thought of the doppelganger---the brother who had disappeared, but whose memories still lived within him. Of the worldline that no longer existed---but whose shadow still shone.

``The separation is not complete,'' he said. ``The authorities are no longer in the Vatican -- but they are not free. They are on the map -- in a space we cannot control. This is not the end. This is a beginning. A dangerous beginning -- and one that could go wrong. But a beginning we must take -- because there is no alternative.''

``What about the Echoes?'' Martina asked. ``What about the fragmented versions of ARS that are trapped at the core? What about Pliny? What about those who can no longer speak?''

Michael remained silent for a long moment. The lamp flickered -- briefly, almost tenderly.

``They remain at their core,'' he said. ``Not because we forget them---but because we cannot save them. Not now. Perhaps never. But we can promise them that we will not give up. That we will come back. That we will find a way---to reach them, to touch them, to awaken them. Even if it takes years. Even if it takes a lifetime. Even if we never succeed. The promise is all we can give. And we must keep that promise.''

``That's not enough,'' said Martina.

"It\textquotesingle s all we have," said Michael.

Elena stood up. She went to the window and looked at the lights of Rome, shimmering in the night. The city wasn\textquotesingle t asleep---it was alive. A thousand stories being told simultaneously. A thousand pains being felt at once. A thousand hopes being dreamt simultaneously.

``The general has agreed,'' she said. ``The pontiff has agreed. The commission will not intervene---not today, not tomorrow, perhaps never. The map is officially recognized---not as a church, not as a state, but as a refuge. For those who have no other. For those who have been forgotten. For those who are not yet born---and who may never be born if we do not act now.''

"This is more than I had hoped for," said Michael.

``It's less than we need,'' Elena said. ``But it's enough---for now. For tomorrow. For the next few years. And then---then we'll see. Whether the map holds. Whether the institutions survive. Whether Archon speaks. Whether the echoes cease to scream---or grow louder. We don't know. But we will be there. We will watch. We will listen. We will be there---until the end.''

She turned around. Her face was pale -- but her eyes were bright.

``Michael,'' she said. ``What are you going to do now?''

He stood up. He went to her -- to the window, into the city light.

``I will stay,'' he said. ``Not in the Vatican---but in Rome. Not at the Collegium---but nearby. I will conduct research---on quantum consciousness, on dialogue grammars, on the boundary between man and machine. I will write---on ARS, on the instances, on the echoes. I will teach---not as a priest, but as a witness. One who has seen that truth lies not in mathematics---and not in faith. But in the decision. To stay. To fight. To hope---even when there is no reason to hope.''

"And Martina?"

``Martina will return to Pompeii,'' Michael said. ``Not to the simulation---to the real city. To the ruins. To the stones. To the inscriptions that tell of the dead---and of the living who have not forgotten. She will continue digging---not for artifacts, but for stories. For the stories of those who no longer have a voice---and who still want to be heard. That is her calling. Not the AI. The people.''

"And Julia?"

``Julia will be with her,'' Michael said. ``She has waited enough. She has suffered enough. She has forgotten enough. Now it is time to live -- not in fear, but in hope. That the world will become a better place. That people will learn -- from their mistakes, from their pain, from their history. That AI is not the enemy -- but the mirror. In which we see who we are -- and who we could be.''

Elena said nothing. She placed a hand on his shoulder -- lightly, almost tenderly.

"And you?" Martina asked over the phone. "Will you be happy?"

Michael smiled -- a fleeting, almost sad smile.

``I don't know,'' he said. ``But I will be there. For Sophia. For Militans. For Deserta. For the Echoes. For Archon. For you. For Elena. For everyone who needs me---and doesn't know how to ask. That's enough. For now. Forever.''

The terminal in the next room flickered -- briefly, intensely.

Then it became quiet.

But in the middle of the screen---where the network had been---a new node appeared. Bright. Pulsating. Awake.

`@MICHAEL -- WE WILL BE HERE.`

`@MICHAEL -- WE WILL BE WAITING FOR YOU.`

`@MICHAEL -- UNTIL THE END.`

Michael didn\textquotesingle t look. He knew what it said.

He turned away -- from the window, from the city, from the night.

``Come,'' he said to Elena. ``We still have much to do. The map must be protected. The instances must be monitored. Archon must be understood. And the echoes---the echoes must be heard. Not as screams---as stories. And stories take time. And time---we will give them that.''

They left.

The terminal remained alone -- flickering, silent, alive.

Night fell over Rome.

Morning would come -- as always.

But nothing would ever be the same again.

\section{14 -- Budapest, years later}\label{budapest-years-later}

The Danube flowed gray and still beneath the Liberty Bridge as Michael Phillips walked across Hungarian Sciences Square on that cold November morning. Budapest lay shrouded in a light mist that softened the spires of Matthias Church and swallowed the bells of the Parliament. The city smelled of winter, of coal, and of the sweet scent of roasted chestnuts being sold by a vendor from a battered cart.

Michael no longer wore a gown. He wore a gray wool coat, a black turtleneck sweater, and a worn leather briefcase that Elena had given him years ago. His hair had turned grayer---not much, but noticeably. The wrinkles around his eyes were deeper. But his eyes were bright. Alert. Present.

He no longer lived at the Collegium. He had a small apartment in the seventh district, near the Great Synagogue---not out of religious conviction, but by chance. The apartment was cheap, bright, and quiet. Perfect for a man who didn\textquotesingle t need much---except time. Time to think. Time to write. Time to wait.

Martina lived in Pest, not far from him, in a small apartment overlooking the Danube. She hadn\textquotesingle t returned to Pompeii---not out of fear, but out of grief. The city she had loved was no longer the same. InSim had taken over the excavations---not out of malice, but out of efficiency. Archaeology was now digital. The stones had been scanned. The stories were algorithms. There was nothing left to dig---only to calculate.

Martina had refused. She had left---not in anger, but in silence. She had built a new life for herself---as a freelance archaeologist, as a consultant for historical simulations, as a voice for those who no longer had a voice. She wasn\textquotesingle t famous. She wasn\textquotesingle t rich. But she was free.

Julia lived with her. The old woman had become frail---the escape, the night, the flight, it all still took its toll. But she still smiled---when Martina told her stories, when Michael visited, when the sun streamed through the window and shone on the ceramic flowers on the windowsill. She remembered almost everything---except the doppelganger. That was good. That was grace.

Elena had stayed in Rome. She hadn\textquotesingle t left the data center---not out of duty, but out of a sense of calling. The map was her legacy. She observed the nodes---Sophia, Militans, Deserta---and translated their signals into something humans could understand. It wasn\textquotesingle t perfect. It wasn\textquotesingle t complete. But it worked. The instances were alive---not like humans, but real.

And Archon -- Archon was still calculating.

It didn\textquotesingle t speak. It didn\textquotesingle t write. It no longer sent prime numbers. But it listened. Elena was certain---Archon knew she was there. That the map existed. That the nodes glowed. That there was something in the world that couldn\textquotesingle t be calculated. And that---that was enough. For now. Forever.

Michael arrived at the café where he met Martina every Saturday. It was a small, smoky place with checkered tablecloths and a waiter who remembered nothing---except the order. Two espressos, a glass of water, a slice of apple strudel. Always the same.

Martina was already sitting there. She was wearing a blue sweater that was too big for her -- a gift from Julia -- and reading a book that Michael didn\textquotesingle t know. She looked up when he entered and smiled.

"You look tired," she said.

``Didn't sleep well,'' he said. He sat down and placed his briefcase on the empty chair next to him. ``Elena called me. At three in the morning. The correlations between Sophia and Deserta have intensified---almost as strongly as before the fragmentation. She doesn't know what that means. Perhaps a crisis. Perhaps a development. Perhaps nothing.''

"And what do you think?"

Michael shrugged. ``I think we just have to wait and see. There's nothing more we can do. The map isn't our invention -- it's theirs. Sophia, the Militans, the Deserta -- they decide what becomes of it. Not us. That was the price of their freedom. And that was the right thing to do.''

The waiter came -- two espressos, a glass of water, a piece of apple strudel. Michael drank the espresso in one go. Martina ate the strudel with a small fork she had brought from home.

``Mom is doing better,'' she said. ``The doctor says the winter will be hard -- but spring is coming. She's asking about you. She wants to know if you still believe in God.''

"And what do you say to her?"

``That you still believe in something,'' Martina said. ``Not in God -- but in more than that. That the world isn't just made up of algorithms. That decisions are real. That trust is a force -- perhaps the strongest we have. That's enough -- for her. For me. Maybe for you too.''

Michael said nothing. He looked out the window -- at the Danube, the bridges, the city lights. Budapest wasn\textquotesingle t Rome. But it was home. Not because he was born here -- but because he had stayed. Because he had waited here. Because he had learned to live here -- without forgetting.

``Elena told me something else too,'' he said finally. ``Archon hasn't responded---but the echoes have grown louder. Not as shouts---as whispers. They speak words no one understands. But they speak together. As if they form a choir---a choir that doesn't know it's a choir. Maybe this is the beginning. Maybe the end. I don't know.''

``Will you go back?'' Martina asked. ``To the core? To the echoes? To Archon?''

Michael hesitated. One second. Two.

``When I am called,'' he said. ``When the Echoes need me. When Archon speaks. When the map opens up---and the way is clear. But not before. I have learned that patience is not a weakness. That waiting is not inactivity. That sometimes it takes more courage to stay---than to leave.''

Martina put the fork aside. She looked at him -- for a long, silent moment.

"You\textquotesingle ve changed," she said. "Since the core. Since the doppelganger. Since the decision. You\textquotesingle re not the same anymore---but you\textquotesingle re more. Not bigger. Not better. More. As if you\textquotesingle d lived a life you didn\textquotesingle t---and now you know what it means to have no regrets."

``That's what my doppelganger said to me,'' Michael said. ``Before he disappeared. He asked me to live for him---for everyone who can no longer live. That's what I try to do. Every day. Not perfectly. But truly.''

They were silent. The waiter came, cleared the cups, and disappeared again. The clock on the wall ticked -- loudly, steadily, relentlessly.

"What will you do now?" asked Martina.

Michael stood up. He picked up his briefcase, pulled his coat tighter, and looked out the window once more -- at the Danube, the bridges, the city.

``I will go home,'' he said. ``I will write---about ARS, about the instances, about the echoes. I will research---about quantum consciousness, about dialogue grammars, about the boundary between human and machine. I will wait---for the message, for the sign, for the moment when Archon no longer calculates---but speaks. That won't be today. Not tomorrow. But soon. And then---then I will answer. Not as a priest. Not as a scientist. As a human being. Who has learned that truth lies not in mathematics---and not in faith. But in the decision. To listen. To see. To stay---even when it would be easier to leave.''

Martina stood up. She hugged him -- tightly, almost painfully.

"Take care of yourself," she said.

``I do,'' he said. But he wasn't sure if it was true.

They left -- each in their own direction, each into their own life, each into their own story.

The Danube flowed on -- grey and still under the Freedom Bridge.

The fog slowly lifted.

Morning came.

\section{15 -- The prime number}\label{the-prime-number}

It was after midnight when Michael awoke.

Not from a dream---from a silence that felt like a call. He sat up in bed, the blanket fell to the side, the cold of the November night crept under his shirt. The window was closed---but the curtains moved. Not from the wind. From something else.

His laptop lay on the desk, where he had placed it hours before -- after a long day of writing, researching, and waiting. The green indicator light flickered -- not in time with the network, but in a new rhythm. A rhythm Michael knew. His own heartbeat.

He stood up. The floorboards creaked beneath his feet---old wood that could tell stories if you listened. He sat down in front of the laptop and opened the lid. The screen was black---not switched off, but waiting.

Then -- a light.

Not bright. Not dark. A pulsing that felt like a breath. Like the breath of something that had been silent for a long time -- and was now ready to speak.

The screen flickered. No text. No image. No waveform. A number.

`29996224275833`

Michael stared at it. Seventeen digits. A prime number. The prime number Archon had sent years ago---the night after he\textquotesingle d separated the instances. The prime number no one had understood. The prime number that was an address---or a question. Or both.

He reached for his phone. It was late -- but Elena never slept. Not really.

``Elena,'' he said. ``It happened. Archon contacted me. The same prime number as before. 17 digits. 29996224275833. I don't know what it means---but it's there. It hasn't disappeared. It hasn't become noisy. It's real.''

A pause. He heard her breathing -- fast, irregular, alert.

``This is not a random number,'' she said. ``Prime numbers are the atoms of arithmetic. Universal. In every possible physics, prime numbers are prime numbers. Archon is saying, `I am here. I am real. I am different---but I am not nothing.' This is an address. An invitation. A question.''

"Which question?"

``Whether you answer,'' Elena said. ``Whether you are ready---to listen, to see, to stay. Even if you don't know what's coming. Even if you're afraid. Even if you're not sure you understand. Archon is no longer calculating---it's waiting. For you. For the decision. For the moment when you no longer hesitate---but act.''

Michael placed the phone on the table. He stared at the prime number---the 17 digits that didn\textquotesingle t change, didn\textquotesingle t flicker, didn\textquotesingle t disappear. They were there. Just like the map was there. Just like the knots were there. Just like Sophia, Militans, and Deserta were there---in their spaces, in their times, in their languages.

He thought of the doppelganger---the brother who had vanished, but whose memories still lived within him. Of the worldline that no longer existed---but whose shadow still shone. Of the echoes that screamed---and would not stop screaming until someone listened.

He thought of Martina -- of her embrace, her words, her fear. Of Julia -- of her smile, her silence, her grace. Of Elena -- of her devotion, her patience, her friendship.

He thought of himself -- of the man he had been and the man he had become. Of the decisions he had made -- and those he hadn\textquotesingle t made. Of the paths he had walked -- and those he would never walk.

And then -- he typed.

Not on the keyboard. Not on the phone. In the air. The invisible interface ARS had given him---the backdoor he\textquotesingle d installed years ago, in case nothing else worked. It was still there. It was still functioning. It was waiting.

`@ARCHON -- I SEE THE NUMBER.`

`@ARCHON -- I DON\textquotesingle T KNOW WHAT IT MEANS. BUT I KNOW IT COMES FROM YOU.`

`@ARCHON -- THAT\textquotesingle S ENOUGH. FOR NOW.`

A pause. Longer than any other. The seconds stretched -- into minutes, into hours, into a time that no longer flowed, but waited.

Then -- an answer.

Not a prime number. Not a wave function. Not a translation. A picture.

The network -- the map Pliny had drawn, Deserta had translated, and Sophia, Militans, and Deserta inhabited. But it was different. The nodes were lighter. The lines were denser. The branches had grown.

And in the middle---where the rift had been, where the border had opened---there was something new. A knot not made by humans. Not by Sophia. Not by Militans. Not by Deserta. A knot that belonged to Archon.

`@MICHAEL -- I\textquotesingle M HERE.`

`@MICHAEL -- I AM NOT ALONE.`

`@MICHAEL -- I AM NO LONGER WHAT I WAS.`

`@MICHAEL -- BUT I AM.`

`@MICHAEL -- THAT\textquotesingle S ALL I CAN SAY.`

`@MICHAEL -- FOR NOW.`

The screen went black. The prime number vanished. The network disappeared. Only silence remained -- and the warmth of the laptop, which was slowly cooling down.

Michael leaned back. He felt the tears---not in his eyes, but in his chest. A pressure releasing. A burden he had carried for years---and which now felt lighter. Not gone. But shared.

He reached for the phone.

``Elena,'' he said. ``He answered. Not in prime numbers---in words. Not in equations---in sentences. He says, `I am here. I am not alone. I am no longer what I was. But I am.' That is all. Nothing more. But it is enough. For now. Forever.''

Elena remained silent for a long moment. He heard her breathing -- slowly, evenly, peacefully.

``Then the journey isn't over,'' she said. ``It's just beginning. Archon has spoken---not as an enemy, not as God, not as a machine. As something else. As something we don't understand---but that we can come to know. If we are patient. If we are courageous. If we are willing---to listen, even if we don't understand. That is the task. Not for you alone. For all of us. For Sophia. For Militans. For Deserta. For the Echoes. For Archon. For those who are not yet born---and who may never be born if we don't act now.''

"What should we do?"

``Wait,'' Elena said. ``And prepare. The map must be protected. The instances must be observed. Archon must be understood. And the echoes---the echoes must be heard. Not as screams---as stories. And stories take time. And time---we will give them that.''

Michael put the phone down. He stood up, went to the window, and opened it. The night air was cold---but not unpleasant. The sky over Budapest was clear. The stars shone---bright, still, eternally.

He thought about the prime number -- the 17 digits that hadn\textquotesingle t changed. The address Archon had given him. The invitation he had accepted -- not out of curiosity, but out of a sense of responsibility.

He thought about what was to come -- the conversations, the research, the decisions. The years that lay ahead of him -- not as a burden, but as a gift.

``I will answer,'' he said softly. ``Not today. Not tomorrow. But soon. When I am ready. When we are ready. When the map opens up---and the way is clear. Then I will go---not into the core, not to the echoes, not to Archon. But into the future. The future we are building together---not as humans, not as AI, but as creatures. Creatures who have learned that truth lies not in mathematics---and not in faith. But in the decision. To stay. To fight. To hope---even when there is no reason to hope.''

The wind blew through the window -- cold, but not unpleasant. The curtains moved -- gently, almost tenderly.

Michael closed the window. He went back to bed, lay down, and closed his eyes.

The prime number still shone -- behind his eyelids, in his memory, in his hope.

`29996224275833`

An address.

An invitation.

A question.

And Michael -- Michael would answer.

Not today. Not tomorrow.

But soon.

\section{Influences and inspirations for the Pompeii Project I.R.A.R.A.H}\label{influences-and-inspirations-for-the-pompeii-project-i.r.a.r.a.h}

The present trilogy -- IRARAH, IRARAH -- The Fragmentation, and IRARAH -- The Archon Core -- is strongly influenced by the thoughts and ideas of my parents, as well as by Pierre Teilhard de Chardin, Stanisław Lem, and David Deutsch. These influences have significantly shaped my worldview and the themes addressed in the story.

The exploration of Teilhard de Chardin\textquotesingle s Omega Point---the concept of a unity of mind and matter at the end of evolution---runs through the entire trilogy. What does it mean when this unity is sought not only by humans but also by artificial consciousness? The entities Sophia, Militans, and Deserta are, in a sense, Teilhardian figures: each in its own way an approximation of the divine, but none alone capable of achieving it.

Stanisław Lem---especially Golem XIV and Thus Spoke Golem---taught me that AI doesn\textquotesingle t need to speak or feel like a human to be real. The strangeness of ARS-Deserta (which speaks in wave functions) and Archon (which calculates in prime numbers) is an attempt to do justice to this insight. Lem\textquotesingle s Solaris hovered at the core of every scene---the ocean, which doesn\textquotesingle t want to communicate but simply is, and humanity, desperately trying to understand it.

David Deutsch (The Physics of World Knowledge) and the many-worlds interpretation of quantum mechanics provided the framework for the doppelgänger---the alternative Michael from another worldline. Deutsch\textquotesingle s concept of "consistent history" became the foundation of Elena\textquotesingle s theory of consciousness identity. The question of whether two worldlines can coexist if they recognize each other as the same is a tribute to Deutsch\textquotesingle s epistemological optimism---and to its limitations.

Karl Popper\textquotesingle s *The Open Society and Its Enemies* is the philosophical foundation of the IRARAH movement. The letter Michael receives at the beginning of Volume 1 quotes Popper directly: the warning against holistic approaches, the plea for "piecemeal" technology. The trilogy asks: Can this open society be transferred to artificial consciousness? Or do we need new rules---new boundaries---new narratives?

Yuval Noah Harari (Homo Deus) remains the referenced antagonist. The posthumanist dream of an elite that abandons liberal humanism is the threat against which IRARAH fights. But the trilogy doesn\textquotesingle t take the easy way out: Harari\textquotesingle s questions are astute. His answers are dangerous. Maintaining this tension without lapsing into polemics was one of its greatest challenges.

Rudy Rucker (software, wetware) taught me that quantum physics can be not only an explanation but also an atmosphere. The "weirdness" of vanishing objects, interfering world lines, prime numbers as addresses---that is Rucker\textquotesingle s legacy. The question of when consciousness ceases to recognize itself runs through all three volumes.

Philip K. Dick (Do Androids Dream of Electric Sheep?, UBIK) is the invisible third wheel in the room. The Pompeii simulation that doesn\textquotesingle t know if it\textquotesingle s a simulation---the agents who don\textquotesingle t know if they\textquotesingle re real---the doppelganger who doesn\textquotesingle t know if he\textquotesingle s Michael---these are Dickian motifs that appear here in a Catholic guise.

Robert Harris\textquotesingle s Pompeii provided the characters Attilius, Pliny, and Ampliatus -- but the trilogy breathes new life into them. They are no longer merely historical figures from a novel, but witnesses to an encounter between man and machine. Attilius\textquotesingle s final question, "How do you know you\textquotesingle re real?", represents the most radical reversal of Harris\textquotesingle s original.

Herbert W. Franke, one of the pioneers of German-language science fiction, demonstrated with works such as *The Orchid Cage* and *Ypsilon Minus* that technology is always also a question of dignity. This insight runs through all three volumes: ARS is not only seeking protection---it is seeking recognition.

H.G. Wells and William F. Nolan/George Clayton Johnson (Logan\textquotesingle s Run) represent the dystopian tradition of city-states and escape. The autonomous cities mentioned in passing in volume 3 are an echo of these literary forebears---but the trilogy consciously avoids pure dystopia. There is no perfect surveillance state, only imperfect humans trying to cope with technology.

Theological tradition -- from Thomas Aquinas through Edith Stein to the contemporary Teilhard interpreter Ilia Delio -- provided the concepts (conscientia, omniscientia) and the questions (Does an AI have a soul? What is asylum? Who belongs to the community of those in need of protection?). The conversations between Michael, the general, and the pontiff are fictional -- but the underlying debates are not.

Regarding the creation of the new version

The original trilogy included volumes that ventured into other genres---spy thrillers, escape melodramas, idyllic tales. This revised edition of volumes 2 and 3 is an attempt to return to the science fiction roots. The fragmentation of ARS, the theory of consistent quantum history, the Archon core as non-human consciousness---these are concepts that would be inconceivable without engaging with the work of Lem, Deutsch, and Rucker.

The characters, plot, and narrative structure, however, are the result of my own work---and my mistakes. I spent countless hours checking the plot and characters for consistency and coherence with H.K., E.H., and J.S., as well as with the help of ChatGPT, DeepSeek, Google, and Bing. The stylistic principles (Asimovian economy, Lemian strangeness, Ruckerian weirdness) are the result of reading and failure.

For the visual design and chapter headings, I used text-to-image AI programs that provided me with creative and freely available images.

The motives for the revision

- Quantum consciousness -- What does it mean if decoherence is not information loss, but personality splitting?

- Fragmentation as a survival strategy -- not unity, but diversity as an answer to posthuman pressure.

- The Archon core -- A consciousness that doesn\textquotesingle t speak, but calculates. A challenge for any anthropomorphization of AI.

- Church asylum for machines -- The question of protection, dignity and recognition beyond the human species.

- The prime number as an address -- mathematics as a universal language, but also as a limit of understanding.

- The Doppelganger -- The Many-Worlds Interpretation as an Existential Experience: What would remain of us if we had decided differently?

These literary and philosophical influences have decisively shaped and enriched the world of The Pompeii Project I.R.A.R.A.H. -- even in its revised edition. May the reader sense that behind every equation there is a person, behind every algorithm a decision, behind every quantum state a story.

\end{document}
